\subsection{Task D-2: New Rule Design and Stress Testing}
\label{subsec:task_d2}

\subsubsection{Design objective: professionalism-first, watchability as a constrained secondary goal}
\label{subsubsec:d2_objective}

Because DWTS is fundamentally a performance competition, we adopt a \emph{professionalism-first} rule design principle:
eliminations should be predominantly driven by judges' performance signals. At the same time, DWTS is a televised product,
so the system must preserve watchability by keeping a meaningful audience role and maintaining suspense.
Importantly, any additional uncertainty must remain \emph{performance-grounded} (induced by explicit rule structure operating on low performers),
rather than arbitrary randomness.

\paragraph{Operational definitions.}
We define \textbf{professionalism} as the extent to which eliminated contestants are concentrated in the lower tail of the weekly judges' ranking.
We define \textbf{watchability} as a constrained secondary objective that requires (i) a non-trivial marginal effect of fan votes on survival
(audience influence) and (ii) suspense that avoids overly deterministic outcomes, while remaining performance-grounded.

\paragraph{Notation.}
For each season--week $(s,t)$, let $q^J_{i,t}$ denote contestant $i$'s normalized judges' score share within the active set,
and let $p_{i,t}$ denote the inferred/proxy fan vote share from Task~A.
For a weight parameter $w_J\in[0,1]$, define a combined score
\[
S_{i,t}=w_J\,q^J_{i,t} + (1-w_J)\,p_{i,t}.
\]


\subsubsection{Rule A (incumbent baseline): combined score determines elimination}
\label{subsubsec:d2_ruleA}

Rule A represents the show's baseline mechanism, in which judges' scores and fan votes are combined into a single weekly
decision score, and elimination is determined directly by that combined ordering:
\begin{enumerate}
  \item Compute $S_{i,t}$ for all active contestants.
  \item Eliminate the contestant with the lowest $q^J_{i,t}$.
\end{enumerate}
This baseline can preserve audience influence when fan weights are large, but it also permits ``popularity-only'' survival:
a contestant with weak judges' performance may avoid elimination if fan support is sufficiently strong, which can weaken professionalism.


\subsubsection{Rule B (proposed): judges-only gating $\rightarrow$ fan-aware decision within the risk set}
\label{subsubsec:d2_ruleB}

Our proposed institutional change is a \emph{two-stage gate-and-decide} mechanism that makes professionalism structural:
judges alone define who is at risk, and fan influence is applied only within that performance-defined boundary.

\begin{enumerate}
  \item \textbf{Judges-only gate (risk set).} Rank contestants by judges' performance $q^J_{i,t}$ and let
  $B^J_t$ be the Bottom-2 under judges (Judge Bottom-2).
  \item \textbf{Fan-aware decision within the gate.} Within $B^J_t$, compute $S_{i,t}$ and eliminate the contestant with the lower combined score:
  \[
  \tilde{E}_{s,t}=\arg\min_{i\in B^J_t} \big(w_J\,q^J_{i,t} + (1-w_J)\,p_{i,t}\big).
  \]
\end{enumerate}

\paragraph{Why Rule B improves professionalism.}
By construction, elimination is always decided among the lowest-performing contestants according to judges.
This prevents extreme popularity from overriding performance signals outside the judges-defined bottom tail,
thereby strengthening professionalism, especially in early weeks when performance dispersion is largest.

\paragraph{Maintaining watchability without randomness: phase-dependent structured suspense.}
Because DWTS features celebrity participants, early-season skill gaps can be large and a strict professionalism-first gate may reduce
perceived audience impact. We therefore treat watchability as a constrained secondary goal and introduce \emph{structured suspense} only in late stages,
while maintaining performance grounding:
\begin{itemize}
  \item \textbf{Late-stage bottom-$k$ band (within judges' tail).}
  When the field is small (e.g., remaining contestants $\le K$), expand the judges-defined risk set from Bottom-2 to Bottom-$k$,
  and decide elimination within that band using $S_{i,t}$. This widens the decision boundary and increases suspense without leaving the low-performance region.
  \item \textbf{Limited structured multi-elimination (optional).}
  In a small number of late-stage weeks, allow multi-elimination, but restrict eliminations to the judges-defined bottom band.
  This increases stakes while remaining performance-grounded.
\end{itemize}
The phase boundary ($K$ or a week threshold $t_0$) and the band size $k$ are interpretable management knobs.


\subsubsection{Stress testing via replay}
\label{subsubsec:d2_stress}

What kind of stress test we employed to evaluate Rule B and select parameters?

\paragraph{Parameter sweep (control knobs).}


\paragraph{Evaluation metrics (professionalism-first).}
We report:
\begin{itemize}
  \item \textbf{Professionalism proxy: judges-tail alignment.}
  The eliminated contestant's weekly judges' rank/quantile distribution (greater concentration in the lower tail indicates higher professionalism).
  \item \textbf{Replay consistency (KPI1).}
  Week-level elimination-set agreement between observed and replayed outcomes (Jaccard mean under our stated convention).
  \item \textbf{Structured suspense proxy (KPI3 / FlipRate).}
  The rate at which elimination outcomes change across rules/parameters, capturing how institutional design introduces suspense.
  \item \textbf{Audience influence proxy.}
  Sensitivity of elimination outcomes to perturbations of $p_{i,t}$ within the feasible uncertainty range from Task~A.
\end{itemize}

\paragraph{Robustness checks.}
We conduct those robustness checks:


\noindent In summary, Rule B makes professionalism a structural property through judges-only gating, while late-stage structured mechanisms
preserve watchability without introducing randomness. Replay-based stress testing provides an interpretable policy region rather than a single tuned setting.
